\chapter{Stochastic Resonance in a Spiking Neural Network Model of the Honeybee Antennal Lobe}
\section{Introduction}

One of the first well estalblished findings in neuroscience with the advent of electrical recordings of neuronal activity, was the variability of neurons' responses. 
That is, when stimulated with the same input, neurons respond with different temporal patterns of spikes, while the overall firing rate remains roughly equal \cite{tomkoNeuronalVariabilityNonstationary1974,mainenReliabilitySpikeTiming1995,tolhurstStatisticalReliabilitySignals1983}.

This trial-to-trial variability is commonly called noise. 
Sources of noise in neurons are many and not all variability in response is due to noise (see \parencite{faisalNoiseNervousSystem2008} for a review). 
Still, due to the physical scale of cells, a stochastic component in neuronal function is inevitable by thermal noise alone. % cite something
It is then to be expected that neural systems would have evolved to tolerate or, more daringly, exploit the presence of noise. 

The idea of studying the effects of noise on the dynamics of a system first appeared in the work of \textcite{benziMechanismStochasticResonance1981,benziStochasticResonanceClimatic1982}, where the term Stochastic Resonance (SR) was first used to describe the observed peak in the power spectrum of the solution of the stochastic differential equation in the form

\begin{equation}
    \odv{x}{t} = x\left(a-x^2\right) + A\cos(\Omega t) + \varepsilon\xi(t)
\end{equation}

where $\xi$ is gaussian white noise, for a value of the noise scaling $\varepsilon \in \left( \varepsilon_1, \varepsilon_2 \right)$.
% just placeholder to understand the trajectory to be followed
Then, in a "sensing system" sense, the (relatively) weak $\xi$ helped to recover the weak forcing frequency $\Omega$, which would have not been detected otherwise. 
In this sense noise may be seen as beneficial in signal processing, neurons of the first part of the sensory systems are particularly noisy signal processors.

